\subsection{Lee--Yang 分支系数的分辨子算术、周期化卷积障碍与零温四态地基 \texorpdfstring{SFT}{SFT}}\label{subsec:conclusion-leyang-resolvent-euler-riesz-field-stratification}
本小节把 Lee--Yang 分支系数三次方程的二次分辨子、其模 \(p\) 分解奇偶律、黄金参数二点 Bernoulli 卷积模 \(2\) 周期化后的卷积算子障碍、真实输入 \(40\) 状态核零温四态地基 SFT 的闭式动力学，以及局部累积量与分支三次数域的线性无交压缩为同一组可直接引用的结论。记
\[
u:=\kappa_*^2,\qquad
b:=B^2,\qquad
K_{\mathrm{br}}:=\QQ(u)=\QQ(b),
\]
并令
\[
L_{\mathrm{LY}}/\QQ
\]
为该 Lee--Yang 分支三次数域的 Galois 闭包，
\[
K_{\mathrm{LY}}:=L_{\mathrm{LY}}^{A_3}=\QQ(\sqrt{-111})
\]
为其唯一二次正规子域。

\begin{theorem}[Lee--Yang 分支系数三次方程的二次分辨子场恰等于 \texorpdfstring{$K_{\mathrm{LY}}$}{K_LY}]\label{thm:conclusion-leyang-dominant-edge-quadratic-resolvent-minus111}
\leanverified{paper\_conclusion\_leyang\_dominant\_edge\_quadratic\_resolvent\_minus111}
设
\[
h(U):=324U^3-540U^2+333U-74,
\]
并记
\[
L_h:=\mathrm{Spl}(h).
\]
则 \(L_h=L_{\mathrm{LY}}\)，且其唯一二次分辨子场满足
\[
\QQ\bigl(\sqrt{\Disc(h)}\bigr)=L_h^{A_3}=\QQ(\sqrt{-111})=K_{\mathrm{LY}}.
\]
\end{theorem}
\begin{proof}
由定理 \ref{thm:xi-terminal-zm-kappa-square-cubic-field-s3}，\(u\) 满足不可约三次方程 \(h(u)=0\)，且
\[
\Disc(h)=-2^6\cdot 3^9\cdot 37=-(2^3\cdot 3^4)^2\cdot 111,
\]
其分裂域伽罗瓦群同构于 \(S_3\)。
另一方面，由定理 \ref{thm:xi-terminal-zm-stokes-leyang-shared-artin-representation}，
\[
\QQ(b)=\QQ(u)=K_{\mathrm{br}}.
\]
因此 \(L_h\) 正是 \(K_{\mathrm{br}}\) 的 Galois 闭包，也即 \(L_{\mathrm{LY}}\)。

对任一不可约 \(S_3\)-三次多项式，其分裂域的唯一二次子域由判别式平方类决定，故
\[
\QQ\bigl(\sqrt{\Disc(h)}\bigr)=L_h^{A_3}.
\]
而 \(\Disc(h)\) 的平方自由部分为 \(-111\)，于是
\[
L_h^{A_3}=\QQ(\sqrt{-111}).
\]
按照本节记号，这个二次域正是 \(K_{\mathrm{LY}}\)。证毕。
\end{proof}

\begin{corollary}[Lee--Yang 分支系数三次方程模 \texorpdfstring{$p$}{p} 的分解型由 Kronecker 符号完全控制其奇偶性]\label{thm:conclusion-leyang-modp-one-root-law-density-thirds}
沿用定理 \ref{thm:conclusion-leyang-dominant-edge-quadratic-resolvent-minus111} 的记号。设 \(p\nmid 2\cdot 3\cdot 37\) 为奇素数。则：
\begin{enumerate}
  \item 若
  \[
  \left(\frac{-111}{p}\right)=-1,
  \]
  则 \(\mathrm{Frob}_p\in \Gal(L_{\mathrm{LY}}/\QQ)\cong S_3\) 为奇置换，故
  \[
  h(U)\bmod p
  \]
  的分解型必为 \((1,2)\)；
  \item 若
  \[
  \left(\frac{-111}{p}\right)=1,
  \]
  则 \(\mathrm{Frob}_p\) 为偶置换，故
  \[
  h(U)\bmod p
  \]
  的分解型只能是 \((1,1,1)\) 或 \((3)\)。
\end{enumerate}
特别地，
\[
\left(\frac{-111}{p}\right)=-1
\]
时不可能出现完全分裂或不可约三次，而
\[
\left(\frac{-111}{p}\right)=1
\]
时不可能出现 \((1,2)\) 型分解。
\end{corollary}
\begin{proof}
由定理 \ref{thm:conclusion-leyang-dominant-edge-quadratic-resolvent-minus111}，
\[
K_{\mathrm{LY}}=L_{\mathrm{LY}}^{A_3}=\QQ(\sqrt{-111}).
\]
因此对 \(p\nmid 2\cdot 3\cdot 37\)，Frobenius 在二次子域上的像满足
\[
\left(\frac{-111}{p}\right)=1
\iff
\Frob_p\in A_3,
\qquad
\left(\frac{-111}{p}\right)=-1
\iff
\Frob_p\in S_3\setminus A_3.
\]
另一方面，三次多项式 \(h(U)\bmod p\) 的分解型由同一 Frobenius 共轭类
\[
\Frob_p\in \Gal(L_{\mathrm{LY}}/\QQ)\cong S_3
\]
决定。
\(S_3\) 的三种循环型与三次多项式在 \(\FF_p\) 上的三种分解型一一对应：恒等元对应 \((1,1,1)\)，换位对应 \((1,2)\)，三循环对应 \((3)\)。
于是奇置换恰对应 \((1,2)\)，而偶置换则只能对应 \((1,1,1)\) 或 \((3)\)。证毕。
\end{proof}

\begin{proposition}[Lee--Yang 分支域的 abelian 阴影与非阿贝尔核心分解]\label{prop:conclusion-leyang-abelian-shadow-nonabelian-core}
\(L_{\mathrm{LY}}/\QQ\) 的最大 abelian 子扩张恰为
\[
\boxed{\;
K_{\mathrm{LY}}=\QQ(\sqrt{-111}).\;}
\]
等价地，任何仅由 abelian Artin 角色可见的 Lee--Yang 分支算术，都压缩为单一二次特征
\[
\left(\frac{-111}{p}\right).
\]
相反，若要区分完全分裂、一根型与不可约型三种 Frobenius 分解，或恢复完整的分支排列，则必须进入
\[
\Gal(L_{\mathrm{LY}}/\QQ)\cong S_3
\]
的非阿贝尔部分。
\end{proposition}
\begin{proof}
由定理 \ref{thm:xi-time-part9c-leyang-cubic-artin-weight1}，
\[
\Gal(L_{\mathrm{LY}}/\QQ)\cong S_3.
\]
而
\[
S_3^{\mathrm{ab}}=S_3/[S_3,S_3]\cong C_2,\qquad [S_3,S_3]=A_3.
\]
因此一切阿贝尔商都必须经由符号映射 \(S_3\twoheadrightarrow C_2\)；其对应固定域正是 \(L_{\mathrm{LY}}^{A_3}\)。
由定理 \ref{thm:conclusion-leyang-dominant-edge-quadratic-resolvent-minus111}，
\[
L_{\mathrm{LY}}^{A_3}=K_{\mathrm{LY}}=\QQ(\sqrt{-111}).
\]
这证明了最大 abelian 子扩张的刻画。

至于最后一条，只需注意：符号角色只能区分换位与偶置换，因而会把恒等元与三循环都压缩到同一取值 \(+1\)。
故完整的三分解型、分支排列与相应非平凡振幅数据均不能由 abelian 阴影单独恢复。证毕。
\end{proof}

\begin{theorem}[Euler--Riesz 二相分裂]\label{thm:conclusion-euler-riesz-two-phase-splitting}
在同一 Fibonacci 频率骨架下，存在两条严格不同的规范通道：
\begin{enumerate}
  \item 标量规范通道 \(\log C_m\) 具有完整的 Euler 型渐近展开：对任意固定 \(R\ge 1\)，
  \[
  \log C_m
  =
  \log(2\pi)
  +
  \sum_{r=1}^{R}
  \frac{B_{2r}}{r(2r-1)}
  \Bigl(\varphi^{-1}+\varphi^{2r-3}\Bigr)
  \Bigl(\frac{\varphi}{2}\Bigr)^{(2r-1)m}
  +
  O\!\Bigl(\Bigl(\frac{\varphi}{2}\Bigr)^{(2R+1)m}\Bigr),
  \]
  因而其全部系数等价地由偶 Bernoulli 数与全部偶值 \(\zeta(2r)\) 唯一决定；
  \item 与之共享同一频率骨架的振荡频谱通道则呈现奇异层：若记
  \[
  \nu:=(\RR\to\TT_2)_\#\mu_{\varphi^{-1}}^{(2)},
  \qquad
  a_m(k):=\frac{|\widehat d_m(k)|}{2^m},
  \]
  则 \(\nu\) 在 \(\TT_2\) 上不具有 \(L^2\) 密度；同时有限尺度能量剖面满足
  \[
  a_m(k)^2
  =
  2^{-m}\prod_{j=1}^{m}
  \left(
  1+\cos\frac{2\pi kF_{j+1}}{F_{m+2}}
  \right),
  \]
  即为 Fibonacci 频率驱动的有限 Riesz 乘积。
\end{enumerate}
因此，本文内部同时内生出一个 even-\(\zeta\)-可读的平滑 Euler 层与一个有限 Riesz 乘积型的奇异振荡层；二者不能压缩到同一个 \(L^2\)-Fourier 描述中。
\end{theorem}
\begin{proof}
第一条正是定理 \ref{thm:fold-bin-gauge-constant-stirling-bernoulli-hierarchy} 的展开式；而由定理 \ref{thm:xi-time-part60aa-gauge-constant-explicit-even-zeta-recovery}，该展开的每一阶系数都可显式反演为 \(\zeta(2r)\)。

第二条中，定理 \ref{thm:xi-bernoulli-two-periodization-no-lp-density} 已证明更强结论
\[
\widehat\nu\notin \ell^r(\ZZ)\qquad (1\le r<\infty),
\]
故特别推出 \(\nu\) 不具有 \(L^2(\TT_2)\) 密度。另一方面，由结论 \ref{con:xi-fold-collision-spectral-trace-parseval} 与定理 \ref{thm:fold-spectrum-factorization}，
\[
a_m(k)^2
=
\frac{|\widehat d_m(k)|^2}{4^m}
=
\prod_{j=1}^{m}\cos^2\!\left(\frac{\pi kF_{j+1}}{F_{m+2}}\right)
=
2^{-m}\prod_{j=1}^{m}
\left(
1+\cos\frac{2\pi kF_{j+1}}{F_{m+2}}
\right),
\]
这正是有限 Riesz 乘积表示。

若两条通道可由同一个 \(L^2\)-Fourier 对象统一承载，则振荡侧在周期化后亦应给出 \(L^2(\TT_2)\) 密度；这与上式所示奇异性相矛盾。故二者在结构上不可压缩为同一个 \(L^2\)-Fourier 描述。证毕。
\end{proof}

\begin{theorem}[周期化测度诱导的卷积算子既非 Hilbert--Schmidt，亦非 trace class]\label{thm:conclusion-periodized-convolution-neither-hs-nor-traceclass}
\leanverified{paper\_conclusion\_periodized\_convolution\_neither\_hs\_nor\_traceclass}
沿用定理 \ref{thm:xi-bernoulli-two-periodization-no-lp-density} 的记号，令
\[
q:=\varphi^{-1},
\qquad
\nu:=\bigl(\RR\to\TT_2=\RR/(2\ZZ)\bigr)_\#\mu_q^{(2)},
\]
并定义卷积算子
\[
\mathcal T_\nu:L^2(\TT_2)\longrightarrow L^2(\TT_2),
\qquad
(\mathcal T_\nu f)(x):=\int_{\TT_2}f(x-y)\,\mathrm{d}\nu(y).
\]
则：
\begin{enumerate}
  \item \(\mathcal T_\nu\) 是有界正规算子；
  \item \(\mathcal T_\nu\) 不是 Hilbert--Schmidt 算子；
  \item \(\mathcal T_\nu\) 不是 trace-class 算子；
  \item 不存在任何 \(K\in L^2(\TT_2\times \TT_2)\) 使
  \[
  (\mathcal T_\nu f)(x)=\int_{\TT_2}K(x,y)f(y)\,\mathrm{d}y
  \]
  在 \(L^2\)-意义下成立。
\end{enumerate}
\end{theorem}
\begin{proof}
对每个 \(m\in\ZZ\)，记
\[
e_m(x):=e^{\pi i m x}\qquad (x\in\TT_2).
\]
则 \((e_m)_{m\in\ZZ}\) 构成 \(L^2(\TT_2)\) 的一组 Fourier 正交基，且
\[
\mathcal T_\nu e_m=\widehat\nu(m)\,e_m.
\]
由 \(\nu\) 为有限测度，得
\[
\sup_{m\in\ZZ}|\widehat\nu(m)|\le \|\nu\|<\infty,
\]
故 \(\mathcal T_\nu\) 有界；又它在一组正交基上可对角化，因此为正规算子。

Hilbert--Schmidt 范数满足
\[
\|\mathcal T_\nu\|_{\mathrm{HS}}^2
=
\sum_{m\in\ZZ}|\widehat\nu(m)|^2.
\]
而定理 \ref{thm:xi-bernoulli-two-periodization-no-lp-density} 已证明
\[
(\widehat\nu(m))_{m\in\ZZ}\notin \ell^r(\ZZ)
\qquad (1\le r<\infty),
\]
尤其不属于 \(\ell^2(\ZZ)\)。故上式发散，\(\mathcal T_\nu\) 不是 Hilbert--Schmidt 算子。

另一方面，任一 trace-class 算子的奇异值序列属于 \(\ell^1\subset \ell^2\)，因此 trace-class 必推出 Hilbert--Schmidt。
既然 \(\mathcal T_\nu\) 不是 Hilbert--Schmidt，它便不可能是 trace class。

最后，若存在 \(K\in L^2(\TT_2\times \TT_2)\) 实现该积分核表示，则 \(\mathcal T_\nu\) 必为 Hilbert--Schmidt，且
\[
\|\mathcal T_\nu\|_{\mathrm{HS}}=\|K\|_{L^2(\TT_2\times\TT_2)}.
\]
这与上面的非 Hilbert--Schmidt 结论矛盾。证毕。
\end{proof}

\begin{theorem}[零温极限四状态地基 \texorpdfstring{SFT}{SFT} 的 \texorpdfstring{$\zeta$}{zeta} 函数、熵与 primitive 轨道渐近具有完全闭式]\label{thm:conclusion-realinput40-zero-temp-fourstate-sft-closed-zeta}
\leanverified{paper\_conclusion\_realinput40\_zero\_temp\_fourstate\_sft\_closed\_zeta}
令
\[
B_\infty=
\begin{pmatrix}
0&1&1&0\\
0&0&1&0\\
0&0&3&1\\
1&0&0&3
\end{pmatrix},
\]
并令 \(\Sigma_{B_\infty}\) 为由 \(B_\infty\) 定义的一侧 SFT，\(\sigma_{B_\infty}\) 为其移位。记
\[
\rho:=\rho(B_\infty),
\qquad
c:=\sqrt{\rho}.
\]
则成立：
\begin{enumerate}
  \item Artin--Mazur \(\zeta\) 函数满足
  \[
  \zeta_{\Sigma_{B_\infty}}(z)
  =
  \frac{1}{\det(I-zB_\infty)}
  =
  \frac{1}{1-6z+9z^2-z^3-z^4};
  \]
  \item 拓扑熵满足
  \[
  h_{\mathrm{top}}(\Sigma_{B_\infty})=\log\rho(B_\infty)=2\log c;
  \]
  \item 由于 \(B_\infty\) 为 primitive 矩阵，存在常数
  \[
  \vartheta\in(0,\rho)
  \]
  使得
  \[
  \#\operatorname{Fix}(\sigma_{B_\infty}^n)
  =
  \operatorname{Tr}(B_\infty^n)
  =
  \rho^n+O(\vartheta^n);
  \]
  \item 若 \(\mathcal O_n\) 表示长度 \(n\) 的 primitive 轨道数，则
  \[
  \mathcal O_n
  =
  \frac{1}{n}\sum_{d\mid n}\mu(d)\operatorname{Tr}\!\bigl(B_\infty^{n/d}\bigr)
  =
  \frac{\rho^n}{n}+O\!\left(\frac{\vartheta^n}{n}\right).
  \]
\end{enumerate}
\end{theorem}
\begin{proof}
对一步 SFT，Artin--Mazur \(\zeta\) 函数的标准公式给出
\[
\zeta_{\Sigma_A}(z)=\frac{1}{\det(I-zA)}.
\]
将 \(A=B_\infty\) 代入，即得第一式。

第二式由 SFT 的标准熵公式
\[
h_{\mathrm{top}}(\Sigma_A)=\log \rho(A)
\]
直接得到；又由零温极限证书
\[
c=\sqrt{\rho(B_\infty)},
\]
故
\[
h_{\mathrm{top}}(\Sigma_{B_\infty})=\log\rho(B_\infty)=2\log c.
\]

再证第三式。由邻接图可见：从状态 \(1\) 可达 \(2,3\)，从 \(2\) 可达 \(3\)，从 \(3\) 可达 \(4\)，从 \(4\) 可达 \(1\)，故图强连通，\(B_\infty\) 不可约。又状态 \(3\) 与 \(4\) 都有自环，因此图周期为 \(1\)，从而 \(B_\infty\) 为 primitive。
由 Perron--Frobenius 定理，\(\rho\) 为单实主特征值，且非 Perron 特征值的模都严格小于 \(\rho\)。取
\[
\vartheta\in
\left(
\max_{j\ge 2}|\lambda_j|,
\rho
\right),
\]
其中 \(\lambda_2,\lambda_3,\lambda_4\) 为其余三个特征值，即得
\[
\operatorname{Tr}(B_\infty^n)=\rho^n+O(\vartheta^n).
\]
而对一步 SFT，
\[
\#\operatorname{Fix}(\sigma_{B_\infty}^n)=\operatorname{Tr}(B_\infty^n),
\]
故第三式成立。

最后，对 primitive 轨道数应用 Möbius 反演：
\[
\mathcal O_n
=
\frac1n\sum_{d\mid n}\mu(d)\operatorname{Tr}\!\bigl(B_\infty^{n/d}\bigr).
\]
其中 \(d=1\) 给出主项 \(\rho^n+O(\vartheta^n)\)；对任意真约数 \(d\ge 2\)，有 \(n/d\le n/2\)，故
\[
\operatorname{Tr}\!\bigl(B_\infty^{n/d}\bigr)=O(\rho^{n/2}).
\]
适当增大 \(\vartheta\) 而仍保持 \(\vartheta<\rho\)，可进一步令 \(\rho^{1/2}<\vartheta\)，从而上述真约数项总和为 \(O(\vartheta^n)\)。于是
\[
\mathcal O_n=\frac{\rho^n}{n}+O\!\left(\frac{\vartheta^n}{n}\right).
\]
证毕。
\end{proof}

\begin{theorem}[项目内部两层不可约算术：\texorpdfstring{$\QQ(\sqrt5)$}{Q(sqrt5)} 与 Lee--Yang 分支三次数域的线性无交]\label{thm:conclusion-local-cumulant-golden-quadratic-leyang-cubic-disjointness}
记
\[
J_{\mathrm{pert}}
:=
\QQ\bigl(P^{(k)}(0):k\ge 1\bigr),
\]
其中 \(P(\theta)=\log\lambda(e^\theta)\) 为真实输入 \(40\) 状态核的输出位势。则有
\[
J_{\mathrm{pert}}=\QQ(\sqrt5),
\qquad
J_{\mathrm{pert}}\cap K_{\mathrm{br}}=\QQ,
\qquad
[J_{\mathrm{pert}}K_{\mathrm{br}}:\QQ]=6.
\]
换言之，局部输出位势的全部累积量生成一个黄金二次层，而 Lee--Yang 分支系数生成一个非正规三次层；二者不存在任何非平凡公共子域。
\end{theorem}
\begin{proof}
由命题 \ref{prop:real-input-40-output-cumulants-closed}，对一切 \(k\ge 1\) 都有
\[
P^{(k)}(0)\in \QQ(\sqrt5),
\]
故 \(J_{\mathrm{pert}}\subseteq \QQ(\sqrt5)\)。另一方面，表
\ref{tab:real_input_40_output_potential_cumulants_closed} 给出
\[
P''(0)=\frac{1791}{200}-\frac{3983}{1000}\sqrt5\notin \QQ,
\]
故 \(\sqrt5\in J_{\mathrm{pert}}\)，于是
\[
J_{\mathrm{pert}}=\QQ(\sqrt5).
\]

再由定理 \ref{thm:conclusion-leyang-dominant-edge-quadratic-resolvent-minus111} 的证明，
\(
K_{\mathrm{br}}
\)
是次数 \(3\) 的非正规三次域。因此
\[
[J_{\mathrm{pert}}\cap K_{\mathrm{br}}:\QQ]
\mid 2,
\qquad
[J_{\mathrm{pert}}\cap K_{\mathrm{br}}:\QQ]
\mid 3,
\]
从而只能有
\[
J_{\mathrm{pert}}\cap K_{\mathrm{br}}=\QQ.
\]
最后由塔式公式
\[
[J_{\mathrm{pert}}K_{\mathrm{br}}:\QQ]
=
\frac{[J_{\mathrm{pert}}:\QQ]\,[K_{\mathrm{br}}:\QQ]}
{[J_{\mathrm{pert}}\cap K_{\mathrm{br}}:\QQ]}
=
\frac{2\cdot 3}{1}
=6.
\]
证毕。
\end{proof}

\begin{corollary}[Lee--Yang 分支坐标不能被任何有限 abelian residue 账本精确承载]\label{cor:conclusion-leyang-branch-coordinate-no-finite-abelian-ledger}
\leanverified{paper\_conclusion\_leyang\_branch\_coordinate\_no\_finite\_abelian\_ledger}
沿用本节记号。不存在有限 abelian 扩张 \(L/\QQ\) 使
\[
u=\kappa_*^2\in L.
\]
因此，一切纯 cyclotomic、Teichm\"uller 或有限 abelian residue 账本，都不可能对 Lee--Yang 分支坐标给出精确封闭表达。
\end{corollary}
\begin{proof}
由定理 \ref{thm:conclusion-local-cumulant-golden-quadratic-leyang-cubic-disjointness}，
\[
K_{\mathrm{br}}=\QQ(u)
\]
是次数 \(3\) 的非正规三次域。若存在有限 abelian 扩张 \(L/\QQ\) 使 \(u\in L\)，则
\[
K_{\mathrm{br}}\subseteq L.
\]
但 abelian 扩张必为 Galois，且其任意中间域对 \(\QQ\) 也仍是 Galois 扩张；于是 \(K_{\mathrm{br}}/\QQ\) 将被迫为正规扩张，这与 \(K_{\mathrm{br}}\) 的非正规性矛盾。故此类 \(L\) 不存在。证毕。
\end{proof}

\begin{theorem}[正则黄金二次层与 Lee--Yang 奇性三次层的最小正规公共载体]\label{thm:conclusion-leyang-global-local-minimal-normal-closure}
沿用
\[
K_{\mathrm{reg}}:=J_{\mathrm{pert}}=\QQ(\sqrt5),
\qquad
K_{\mathrm{sing}}:=K_{\mathrm{br}}=\QQ(u),
\]
并记 \(H_{\mathrm{sing}}\) 为 \(K_{\mathrm{sing}}\) 的 Galois 闭包，设
\[
\Omega:=K_{\mathrm{reg}}H_{\mathrm{sing}}.
\]
则
\[
[K_{\mathrm{reg}}K_{\mathrm{sing}}:\QQ]=6,
\qquad
[\Omega:\QQ]=12,
\qquad
\Gal(\Omega/\QQ)\cong C_2\times S_3.
\]
并且 \(\Omega\) 正是 \(K_{\mathrm{reg}}K_{\mathrm{sing}}\) 的最小正规闭包。
\end{theorem}
\begin{proof}
由定理 \ref{thm:conclusion-local-cumulant-golden-quadratic-leyang-cubic-disjointness}，
\[
K_{\mathrm{reg}}\cap K_{\mathrm{sing}}=\QQ,
\qquad
[K_{\mathrm{reg}}K_{\mathrm{sing}}:\QQ]=2\cdot 3=6.
\]
再由定理 \ref{thm:conclusion-leyang-dominant-edge-quadratic-resolvent-minus111} 与命题 \ref{prop:conclusion-leyang-abelian-shadow-nonabelian-core}，\(H_{\mathrm{sing}}/\QQ\) 的 Galois 群为 \(S_3\)，且其唯一二次子域为
\[
K_{\mathrm{LY}}=\QQ(\sqrt{-111}).
\]
由于
\[
K_{\mathrm{reg}}=\QQ(\sqrt5)\neq \QQ(\sqrt{-111})=K_{\mathrm{LY}},
\]
故
\[
K_{\mathrm{reg}}\cap H_{\mathrm{sing}}=\QQ.
\]
而 \(K_{\mathrm{reg}}/\QQ\) 与 \(H_{\mathrm{sing}}/\QQ\) 都是 Galois 扩张，因此
\[
[\Omega:\QQ]=[K_{\mathrm{reg}}:\QQ]\,[H_{\mathrm{sing}}:\QQ]=2\cdot 6=12,
\]
且
\[
\Gal(\Omega/\QQ)\cong \Gal(K_{\mathrm{reg}}/\QQ)\times \Gal(H_{\mathrm{sing}}/\QQ)\cong C_2\times S_3.
\]
由于 \(H_{\mathrm{sing}}\) 本身就是 \(K_{\mathrm{sing}}\) 的最小正规闭包，而 \(K_{\mathrm{reg}}\) 已经正规，故 \(\Omega\) 正是 \(K_{\mathrm{reg}}K_{\mathrm{sing}}\) 的最小正规闭包。证毕。
\end{proof}

\begin{corollary}[最大全阿贝尔阴影与被迫出现的第三二次子域]\label{cor:conclusion-leyang-global-local-maximal-abelian-shadow}
\(\Omega/\QQ\) 的极大阿贝尔子扩张恰为
\[
\Omega^{\mathrm{ab}}=\QQ(\sqrt5,\sqrt{-111}),
\]
因此它是一个双二次域，其三个且仅有三个二次子域为
\[
\QQ(\sqrt5),\qquad
\QQ(\sqrt{-111}),\qquad
\QQ(\sqrt{-555}).
\]
\end{corollary}
\begin{proof}
由定理 \ref{thm:conclusion-leyang-global-local-minimal-normal-closure}，
\[
\Gal(\Omega/\QQ)\cong C_2\times S_3.
\]
其交换化为
\[
(C_2\times S_3)^{\mathrm{ab}}\cong C_2\times C_2,
\]
而交换子群正是 \(A_3\subset S_3\)。故极大阿贝尔子扩张为固定域 \(\Omega^{A_3}\)，其次数为 \(4\)。

另一方面，\(\Omega\) 同时包含 \(K_{\mathrm{reg}}=\QQ(\sqrt5)\) 与 \(H_{\mathrm{sing}}\) 的唯一二次子域 \(\QQ(\sqrt{-111})\)，因而
\[
\QQ(\sqrt5,\sqrt{-111})\subseteq \Omega^{A_3}.
\]
左端已经是次数 \(4\) 的双二次域，故必有
\[
\Omega^{\mathrm{ab}}=\Omega^{A_3}=\QQ(\sqrt5,\sqrt{-111}).
\]
双二次域的第三个二次子域由平方类乘积给出，即
\[
\QQ(\sqrt5)\cdot \QQ(\sqrt{-111})=\QQ(\sqrt{-555}).
\]
证毕。
\end{proof}

\begin{theorem}[阿贝尔化必然抹去 Lee--Yang 奇性分支的三张 sheet]\label{thm:conclusion-leyang-global-local-abelianization-kills-sheets}\leanverified{paper\_conclusion\_leyang\_global\_local\_abelianization\_kills\_sheets}
对任意阿贝尔群 \(A\) 与任意群同态
\[
\rho:\Gal(\Omega/\QQ)\longrightarrow A,
\]
\(\rho\) 都必湮灭 \(A_3\subset S_3\)。从而，一切阿贝尔可观测量都无法区分 \(f(T)=324T^3-540T^2+333T-74\) 的三根所对应的三张奇性 sheet。等价地，所有阿贝尔观测都只能下降到结论 \ref{cor:conclusion-leyang-global-local-maximal-abelian-shadow} 的双二次阴影，而绝不可能恢复 \(u\) 本身。
\end{theorem}
\begin{proof}
任意到阿贝尔群的群同态都必经交换化，故必湮灭交换子群
\[
[\Gal(\Omega/\QQ),\Gal(\Omega/\QQ)]=A_3.
\]
但在 \(H_{\mathrm{sing}}\) 中，\(A_3\) 正是循环置换三根的三循环子群。因此经由任意阿贝尔观测，三张奇性 sheet 必落入同一不可区分轨道。

若 \(u\) 可由某个阿贝尔观测恢复，则必有
\[
u\in \Omega^{\mathrm{ab}}.
\]
然而推论 \ref{cor:conclusion-leyang-branch-coordinate-no-finite-abelian-ledger} 已证明：\(u\) 不属于任何有限阿贝尔扩张。由于 \(\Omega^{\mathrm{ab}}/\QQ\) 是有限阿贝尔扩张，矛盾。故 \(u\notin \Omega^{\mathrm{ab}}\)，阿贝尔观测不可能恢复 \(u\)。证毕。
\end{proof}

\begin{theorem}[Lee--Yang 奇性三次数域在 \texorpdfstring{$5$}{5} 处的无分歧循环化]\label{thm:conclusion-leyang-global-local-q5-unramified-cubic}
\leanverified{paper\_conclusion\_leyang\_global\_local\_q5\_unramified\_cubic}
记
\[
K_{\mathrm{sing},5}:=K_{\mathrm{sing}}\otimes_{\QQ}\QQ_5.
\]
则 \(K_{\mathrm{sing},5}\) 是 \(\QQ_5\) 的唯一三次无分歧扩张。特别地，
\[
[K_{\mathrm{sing},5}:\QQ_5]=3,
\qquad
e(K_{\mathrm{sing},5}/\QQ_5)=1,
\qquad
f(K_{\mathrm{sing},5}/\QQ_5)=3,
\]
并且
\[
\Gal(K_{\mathrm{sing},5}/\QQ_5)\cong C_3.
\]
\end{theorem}
\begin{proof}
把
\[
f(T)=324T^3-540T^2+333T-74
\]
模 \(5\) 约化，得到
\[
\bar f(T)\equiv 4T^3+3T+1.
\]
直接代入 \(T=0,1,2,3,4\) 可见 \(\bar f\) 在 \(\FF_5\) 中无根。由于 \(\deg \bar f=3\)，故 \(\bar f\) 在 \(\FF_5[T]\) 中不可约。

另一方面，由定理 \ref{thm:conclusion-leyang-dominant-edge-quadratic-resolvent-minus111} 的证明，\(f\) 的判别式为
\[
\Disc(f)=-2^6\cdot 3^9\cdot 37,
\]
故 \(5\nmid \Disc(f)\)。因此 \(5\) 在 \(K_{\mathrm{sing}}\) 中不分歧；又因 \(\bar f\) 不可约，\(5\) 对应唯一素理想，且剩余次数为 \(3\)。于是
\[
e=1,\qquad f=3.
\]
对局部域而言，无分歧三次扩张必为 Galois 且其 Galois 群循环，故
\[
\Gal(K_{\mathrm{sing},5}/\QQ_5)\cong C_3.
\]
证毕。
\end{proof}

\begin{corollary}[正则--奇性联合 \texorpdfstring{$5$}{5}-进载体是 \texorpdfstring{$(e,f)=(2,3)$}{(e,f)=(2,3)} 的循环六次扩张]\label{cor:conclusion-leyang-global-local-q5-composite-c6}
记
\[
L_5:=\QQ_5(\sqrt5,u).
\]
则
\[
[L_5:\QQ_5]=6,
\qquad
e(L_5/\QQ_5)=2,
\qquad
f(L_5/\QQ_5)=3,
\]
并且
\[
\Gal(L_5/\QQ_5)\cong C_6.
\]
\end{corollary}
\begin{proof}
由推论 \ref{cor:conclusion-fibadic-golden-extension-even-ramification-necessity}，
\[
\QQ_5(\sqrt5)/\QQ_5
\]
是二次全分歧扩张。由定理 \ref{thm:conclusion-leyang-global-local-q5-unramified-cubic}，
\[
K_{\mathrm{sing},5}/\QQ_5
\]
是三次无分歧扩张。无分歧扩张与全分歧扩张交域必为基域，故
\[
\QQ_5(\sqrt5)\cap K_{\mathrm{sing},5}=\QQ_5.
\]
于是
\[
[L_5:\QQ_5]=2\cdot 3=6,
\qquad
e(L_5/\QQ_5)=2,
\qquad
f(L_5/\QQ_5)=3.
\]
两侧子扩张都为 Galois，且 Galois 群分别为 \(C_2\) 与 \(C_3\)；因此合成扩张也为 Galois，且
\[
\Gal(L_5/\QQ_5)\cong C_2\times C_3\cong C_6.
\]
证毕。
\end{proof}

\begin{theorem}[全局非阿贝尔性在 \texorpdfstring{$5$}{5}-进纤维上完全隐没]\label{thm:conclusion-leyang-global-local-nonabelianity-disappears-at-q5}
设 \(\mathfrak P\) 为 \(\Omega\) 上位于 \(5\) 之上的素理想。则其分解群满足
\[
D_{\mathfrak P}\cong C_6.
\]
特别地，
\[
\Gal(\Omega/\QQ)\cong C_2\times S_3
\]
是非阿贝尔群，而所有 \(\QQ_5\)-局部完成所见到的分解群却都是阿贝尔的。故 Lee--Yang 奇性分支的非阿贝尔性并不诞生于 \(5\)-进局部纤维，而只诞生于全局拼接。
\end{theorem}
\begin{proof}
由定理 \ref{thm:conclusion-leyang-global-local-q5-unramified-cubic}，\(K_{\mathrm{sing},5}/\QQ_5\) 已经是 Galois 的循环三次扩张。因此 \(H_{\mathrm{sing}}\) 在 \(5\) 之上的局部完成不会比 \(K_{\mathrm{sing},5}\) 更大；换言之，\(H_{\mathrm{sing}}\) 的任一 \(5\)-进完成都同构于 \(K_{\mathrm{sing},5}\)。

再与 \(\QQ_5(\sqrt5)\) 合成，便得到结论 \ref{cor:conclusion-leyang-global-local-q5-composite-c6} 中的 \(L_5\)。因此
\[
\Omega_{\mathfrak P}\cong L_5.
\]
分解群与局部 Galois 群同构，故
\[
D_{\mathfrak P}\cong \Gal(\Omega_{\mathfrak P}/\QQ_5)\cong \Gal(L_5/\QQ_5)\cong C_6.
\]
这说明 \(S_3\)-型非交换性在 \(5\)-进纤维上完全退化为循环六次局部单调性。证毕。
\end{proof}

\begin{theorem}[Lee--Yang 振幅常数的 \texorpdfstring{$5$}{5}-进提升只增加剩余次数而不增加分歧指数]\label{thm:conclusion-leyang-global-local-amplitude-b-q5-lift}
设 \(B\) 为 Lee--Yang 振幅常数，并记
\[
L_B:=\QQ(B).
\]
则 \(B\) 的最小多项式可取为
\[
h(X)=162X^6+1593X^4+1998X^2+1184.
\]
并且
\[
L_{B,5}:=L_B\otimes_{\QQ}\QQ_5
\]
是 \(\QQ_5\) 的唯一六次无分歧扩张；因此
\[
[L_{B,5}:\QQ_5]=6,
\qquad
e(L_{B,5}/\QQ_5)=1,
\qquad
f(L_{B,5}/\QQ_5)=6.
\]
进一步，
\[
[\QQ(\sqrt5,B):\QQ]=12,
\]
且其 \(\QQ_5\)-完成满足
\[
[\QQ_5(\sqrt5,B):\QQ_5]=12,
\qquad
e=2,
\qquad
f=6,
\]
并有
\[
\Gal(\QQ_5(\sqrt5,B)/\QQ_5)\cong C_2\times C_6.
\]
\end{theorem}
\begin{proof}
由定理 \ref{thm:xi-terminal-zm-stokes-amplitude-degree6-minpoly}，
\[
[\QQ(B):\QQ]=6,
\]
且 \(B\) 的极小多项式正是
\[
h(X)=162X^6+1593X^4+1998X^2+1184.
\]
把 \(h\) 模 \(5\) 约化，得
\[
\bar h(X)\equiv 2X^6-2X^4-2X^2-1.
\]
直接在 \(\FF_5[X]\) 中计算可得
\[
\gcd(\bar h,X^5-X)=\gcd(\bar h,X^{25}-X)=\gcd(\bar h,X^{125}-X)=1.
\]
因此 \(\bar h\) 没有 \(1\)、\(2\)、\(3\) 次因子；而 \(\deg \bar h=6\)，遂 \(\bar h\) 在 \(\FF_5[X]\) 中不可约。

再由定理 \ref{thm:xi-leyang-amplitude-B-sextic-biquadratic-shadow}，
\[
\Disc(h)=-2^{16}\cdot 3^{22}\cdot 11^4\cdot 37^3\cdot 109^4,
\]
故 \(5\nmid \Disc(h)\)。于是 \(L_{B,5}/\QQ_5\) 是六次无分歧扩张，从而必为唯一六次无分歧扩张。

另一方面，\(\QQ_5(\sqrt5)/\QQ_5\) 是二次全分歧扩张，而 \(L_{B,5}/\QQ_5\) 无分歧，因此 \(L_{B,5}\) 不可能包含 \(\QQ_5(\sqrt5)\)。于是
\[
\QQ(B)\cap \QQ(\sqrt5)=\QQ,
\]
从而
\[
[\QQ(\sqrt5,B):\QQ]=2\cdot 6=12.
\]
局部上，无分歧六次扩张与全分歧二次扩张交域平凡，故
\[
[\QQ_5(\sqrt5,B):\QQ_5]=12,
\qquad
e=2,
\qquad
f=6.
\]
两侧局部子扩张都为 Galois，因此其合成扩张也为 Galois，且 Galois 群为
\[
C_2\times C_6.
\]
证毕。
\end{proof}

\begin{theorem}[两切片最小完备阈值与离开有限阿贝尔算术账本的阈值重合]\label{thm:conclusion-leyang-global-local-two-slice-threshold}
在单原子差分通道
\[
c^{(j)}=m\tau_j^E
\qquad (j=1,2)
\]
中，单切片永远不足以唯一恢复参数对 \((m,E)\)，而任意两条不同切片都足以通过
\[
E=\frac{\log(c^{(2)}/c^{(1)})}{\log(\tau_2/\tau_1)},
\qquad
m=c^{(1)}\tau_1^{-E}
\]
精确恢复 \((m,E)\)。当指数参数特化为 Lee--Yang 奇性坐标
\[
E=u=\kappa_\star^2
\]
时，最小解析完备阈值与离开有限阿贝尔算术账本的阈值严格重合：达到“两切片”后即可精确恢复 \(u\)，而所恢复出的 \(u\) 立刻不属于任何有限阿贝尔扩张。
\end{theorem}
\begin{proof}
前文已经证明：单切片不足以唯一恢复 \((m,E)\)，而任意两条不同切片则给出上述显式重构公式。把 \(E\) 特化为 \(u\) 后，便得到
\[
u=\frac{\log(c^{(2)}/c^{(1)})}{\log(\tau_2/\tau_1)}.
\]
因此，两切片已经足以精确恢复奇性坐标本身。

另一方面，由推论 \ref{cor:conclusion-leyang-branch-coordinate-no-finite-abelian-ledger}，\(u\) 不属于任何有限阿贝尔扩张。故一旦跨过“恰好两切片”的最小解析完备阈值，所得参数立即越出全部有限阿贝尔 residue 账本。证毕。
\end{proof}

\begin{theorem}[单切片参数纤维的乘法轨道结构与双切片刚性]\label{thm:conclusion-residue-slice-single-orbit-double-slice-rigidity}
\leanverified{paper\_conclusion\_residue\_slice\_single\_orbit\_double\_slice\_rigidity}
设 \(\tau>0\) 为固定切片参数，并令
\[
c=m\tau^E,
\qquad
m>0,
\qquad
E\in\RR.
\]
定义单切片参数纤维
\[
\mathcal F_{\tau}(c):=\{(m,E)\in\RR_{>0}\times\RR:\ m\tau^E=c\}.
\]
则成立：
\begin{enumerate}
  \item 若 \(\tau\neq 1\)，则对任意一组特解 \((m_0,E_0)\in \mathcal F_{\tau}(c)\)，有
        \[
        \mathcal F_{\tau}(c)
        =
        \{(m_0\tau^{-t},\,E_0+t):\ t\in\RR\};
        \]
  \item 若 \(\tau=1\)，则
        \[
        \mathcal F_{1}(c)=\{(c,E):\ E\in\RR\};
        \]
  \item 若 \(\tau_1\neq \tau_2\) 且
        \[
        c_j=m\tau_j^E
        \qquad (j=1,2),
        \]
        则
        \[
        E=\frac{\log(c_2/c_1)}{\log(\tau_2/\tau_1)},
        \qquad
        m=c_1\tau_1^{-E},
        \]
        因而 \((m,E)\) 由两条不同切片唯一恢复。
\end{enumerate}
特别地，单切片只给出一条一维乘法轨道，而双切片恰把该轨道塌缩为点。
\end{theorem}
\begin{proof}
若 \(\tau\neq 1\)，并取 \((m_0,E_0)\in\mathcal F_{\tau}(c)\)，则对任意 \(t\in\RR\) 定义
\[
m_t:=m_0\tau^{-t},
\qquad
E_t:=E_0+t.
\]
直接计算得
\[
m_t\tau^{E_t}=m_0\tau^{-t}\tau^{E_0+t}=m_0\tau^{E_0}=c,
\]
故 \((m_t,E_t)\in \mathcal F_{\tau}(c)\)。反之，任取 \((m,E)\in \mathcal F_{\tau}(c)\)，令 \(t:=E-E_0\)，则
\[
m\tau^{E}=m_0\tau^{E_0}
\]
推出
\[
m=m_0\tau^{E_0-E}=m_0\tau^{-t},
\]
从而 \((m,E)\) 恰可写成上述形式。

若 \(\tau=1\)，则约束方程退化为 \(m=c\)，而 \(E\) 自由，故得到第二条。

最后，若 \(\tau_1\neq \tau_2\)，则对两式相除得
\[
\frac{c_2}{c_1}=\left(\frac{\tau_2}{\tau_1}\right)^E.
\]
由于 \(\tau_2/\tau_1\neq 1\)，故
\[
E=\frac{\log(c_2/c_1)}{\log(\tau_2/\tau_1)}.
\]
再回代 \(c_1=m\tau_1^E\) 即得
\[
m=c_1\tau_1^{-E}.
\]
证毕。
\end{proof}

\begin{theorem}[有限 cyclotomic 补丁不能替代缺失的第二切片]\label{thm:conclusion-residue-cyclotomic-patch-cannot-replace-second-slice}
\leanverified{paper\_conclusion\_residue\_cyclotomic\_patch\_cannot\_replace\_second\_slice}
设
\[
G^{(N)}(z)=G^{\mathrm{core}}(z)+\sum_{\nu=1}^{N}\frac{2z^{\ell_\nu}}{1-z^{M_\nu}},
\qquad
M_\nu\ge 2,
\]
并沿用定理 \ref{thm:conclusion-finite-cyclotomic-surgery-positive-exponent-transparency} 的指数参数
\[
\theta^{\mathrm{core}},
\qquad
\theta^{(N)}.
\]
则恒有
\[
\theta^{(N)}=\max\{0,\theta^{\mathrm{core}}\}.
\]
特别地，任意有限 cyclotomic 补丁都不改变正 twisted 指数的阈值结构；因此它也不改变定理 \ref{thm:conclusion-residue-slice-single-orbit-double-slice-rigidity} 中的最小切片完备性。换言之，若只给定一条切片上的数据
\[
c=m\tau^E,
\]
则在加入任意有限 cyclotomic 补丁之后，参数纤维 \(\mathcal F_{\tau}(c)\) 仍然是一维的；有限 cyclotomic 手术不能充当缺失的第二切片。
\end{theorem}
\begin{proof}
第一条正是定理 \ref{thm:conclusion-finite-cyclotomic-surgery-positive-exponent-transparency} 的内容。该定理说明：有限 cyclotomic 补丁只会在单位圆上添加奇点，而不会制造任何新的正 twisted 指数门槛。

因此，在恢复 \((m,E)\) 的问题上，有限 cyclotomic 补丁并不会产生额外的独立切片参数；单切片所留下的自由度仍完全由方程
\[
m\tau^E=c
\]
控制。于是定理 \ref{thm:conclusion-residue-slice-single-orbit-double-slice-rigidity} 立即给出：单切片纤维仍为一维乘法轨道，而唯一使其塌缩的最小数据仍是第二条不同切片。证毕。
\end{proof}

\begin{theorem}[prime-power 完成化塔在 \texorpdfstring{\(\FF_p\)}{Fp} 上塌缩为单生成 Frobenius 代数]\label{thm:conclusion-completed-primepower-frobenius-single-generator-collapse}
\leanverified{paper\_conclusion\_completed\_primepower\_frobenius\_single\_generator\_collapse}
对任意素数 \(p\)，定义
\[
\mathcal A_p:=
\FF_p\bigl[\widehat a_p(s),\widehat a_{p^2}(s),\widehat a_{p^3}(s),\dots\bigr]
\subseteq \FF_p[s].
\]
则有
\[
\widehat a_p(s)\equiv s^p\pmod p,
\qquad
\widehat a_{p^k}(s)\equiv s^{p^k}\pmod p
\qquad (k\ge 1),
\]
从而
\[
\mathcal A_p=\FF_p[s^p].
\]
因此，整条 prime-power 完成化塔在 \(\bmod p\) 上并不生成真正多参数的坐标代数，而只是单生成元 \(s^p\) 的 Frobenius 迭代轨道。
\end{theorem}
\begin{proof}
由定理 \ref{thm:conclusion-completed-prime-length-frobenius-collapse}，
\[
\widehat a_p(s)\equiv s^p\pmod p.
\]
再由定理 \ref{thm:conclusion-completed-primepower-dwork-frobenius-tower}，
\[
\widehat a_{p^k}(s)\equiv s^{p^k}\pmod p
\qquad (k\ge 1).
\]
故每个 \(\widehat a_{p^k}(s)\) 都属于 \(\FF_p[s^p]\)，从而
\[
\mathcal A_p\subseteq \FF_p[s^p].
\]
另一方面，\(\widehat a_p(s)\equiv s^p\pmod p\) 已表明 \(s^p\in \mathcal A_p\)，因而
\[
\FF_p[s^p]\subseteq \mathcal A_p.
\]
两边合并即得
\[
\mathcal A_p=\FF_p[s^p].
\]
证毕。
\end{proof}

\begin{corollary}[regular germ 与 singular branch 不能被同一个有限 Abelian residue ledger 精确承载]\label{cor:conclusion-regular-singular-no-finite-abelian-ledger}
\leanverified{paper\_conclusion\_regular\_singular\_no\_finite\_abelian\_ledger}
设
\[
K_{\mathrm{reg}}:=\QQ(\sqrt5),
\qquad
K_{\mathrm{sing}}:=\QQ(u_\ast),
\qquad
u_\ast=\kappa_\star^2.
\]
则不存在有限 Abelian 扩张 \(L/\QQ\) 使
\[
K_{\mathrm{reg}}K_{\mathrm{sing}}\subseteq L.
\]
特别地，不存在任何纯 cyclotomic、纯 Teichm\"uller 或纯有限 Abelian residue ledger，能够同时精确承载 regular side 的局部累积量与 singular side 的 Lee--Yang 分支坐标。
\end{corollary}
\begin{proof}
反设存在有限 Abelian 扩张 \(L/\QQ\) 使
\[
K_{\mathrm{reg}}K_{\mathrm{sing}}\subseteq L.
\]
由于有限 Abelian 扩张必为 Galois，且其任意中间扩张的正规闭包仍是 Abelian，故 \(K_{\mathrm{reg}}K_{\mathrm{sing}}\) 的最小正规闭包也必须是 Abelian。

然而由定理 \ref{thm:conclusion-leyang-global-local-minimal-normal-closure}，
\[
\Gal(\Omega/\QQ)\cong C_2\times S_3
\]
正是 \(K_{\mathrm{reg}}K_{\mathrm{sing}}\) 的最小正规闭包的 Galois 群；该群非 Abelian，矛盾。故此类 \(L\) 不存在。证毕。
\end{proof}

\begin{theorem}[completion-faithful residue 证书的三层不可省略分裂]\label{thm:conclusion-completion-faithful-certificate-three-register-split}
设 \(\mathcal C\) 为一类 completion-side 证书，目标是在同一单原子 residue 通道上忠实恢复四元组
\[
(\ell,m,E,u_\ast).
\]
则 \(\mathcal C\) 至少必须同时暴露下列三类不可省略的读出层：
\[
\mathfrak R_{\mathrm{CRT}}
\quad\text{用于恢复离散支撑 }\ell;
\]
\[
\mathfrak R_{\mathrm{slice}}
\quad\text{用于以两条不同切片恢复 }(m,E);
\]
\[
\mathfrak R_{\mathrm{branch}}
\quad\text{用于承载 }u_\ast\text{ 的非 Abelian 奇性分支}.
\]
更精确地，任何省去其中任意一层的方案都不可能 faithful。
\end{theorem}
\begin{proof}
首先，由定理 \ref{thm:conclusion-atomic-surgery-finite-congruence-tomography}，有限互素模层析通过中国剩余定理唯一恢复离散支撑 \(\ell\)。故若省去 \(\mathfrak R_{\mathrm{CRT}}\)，则 \(\ell\) 无法被忠实恢复。

其次，由定理 \ref{thm:conclusion-residue-slice-single-orbit-double-slice-rigidity}，单切片方程
\[
m\tau^E=c
\]
总留下一个一维参数轨道；而定理 \ref{thm:conclusion-residue-cyclotomic-patch-cannot-replace-second-slice} 进一步说明，任何有限 cyclotomic 补丁都不能替代缺失的第二切片。故若省去 \(\mathfrak R_{\mathrm{slice}}\)，则 \((m,E)\) 不能被唯一恢复。

最后，由推论 \ref{cor:conclusion-regular-singular-no-finite-abelian-ledger}，\(u_\ast\) 不能与 regular side 的算术数据一起落入任何有限 Abelian residue ledger。故若省去 \(\mathfrak R_{\mathrm{branch}}\)，则 singular branch 信息必被压到 Abelian 阴影中，从而不能 faithful 地恢复 \(u_\ast\)。

因此，对同一通道的 faithful 恢复必须同时保留这三层。证毕。
\end{proof}

\begin{corollary}[若再要求乘法 bookkeeping faithful，则必须外置第四实现层]\label{cor:conclusion-completion-faithful-certificate-fourth-implementation-layer}
在定理 \ref{thm:conclusion-completion-faithful-certificate-three-register-split} 的语义下，若还要求证书机制对 residue mass 的乘法结构保持 faithful bookkeeping，则除
\[
\mathfrak R_{\mathrm{CRT}}
\oplus
\mathfrak R_{\mathrm{slice}}
\oplus
\mathfrak R_{\mathrm{branch}}
\]
之外，还必须额外外置一层
\[
\mathfrak R_{\mathrm{impl}}
\]
作为实现层寄存器。原因在于
\[
\cdim^{\mathrm{impl}}(\NN_\times)=\frac12,
\qquad
\cdim^{\mathrm{hom}}_{+}(\NN_\times)=+\infty.
\]
因此，residue mass 作为集合可以被单寄存器实现，但作为乘法幺半群却不能被任何有限维加法账本忠实同态化。
\end{corollary}
\begin{proof}
定理 \ref{thm:conclusion-completion-faithful-certificate-three-register-split} 已经说明：若要 faithful 地恢复 \((\ell,m,E,u_\ast)\)，至少必须保留
\(
\mathfrak R_{\mathrm{CRT}},
\mathfrak R_{\mathrm{slice}},
\mathfrak R_{\mathrm{branch}}
\)
三层。

若进一步要求 residue mass 的乘法 bookkeeping 也被 faithful 保留，则由定理 \ref{thm:conclusion-implementation-compression-hom-linearization-split}，
\[
\cdim^{\mathrm{impl}}(\NN_\times)=\frac12
\]
说明 \(\NN_\times\) 在实现层可以被单寄存器注入；但
\[
\cdim^{\mathrm{hom}}_{+}(\NN_\times)=+\infty
\]
说明它不可能被任何有限维加法寄存器忠实同态化。故除前三层外，还必须额外外置一层真正负责乘法实现的 \(\mathfrak R_{\mathrm{impl}}\)。证毕。
\end{proof}

\begin{conjecture}[completion-faithful \texorpdfstring{\(\pi\)}{pi}-公式的 residue--slice--branch--implementation 四层分裂]\label{conj:conclusion-piformula-residue-slice-branch-implementation-split}
任一若要在本文的 completion 语义中不仅逼近数值 \(\pi\)，而且 faithful 地承载 residue/branch 结构的公式，都至少应当分裂为以下四层：
\[
\mathfrak R_{\mathrm{res}}:=\mathfrak R_{\mathrm{CRT}},
\qquad
\mathfrak R_{\mathrm{amp}}:=\mathfrak R_{\mathrm{slice}},
\qquad
\mathfrak R_{\mathrm{branch}},
\qquad
\mathfrak R_{\mathrm{impl}}.
\]
其中：
\begin{enumerate}
  \item \(\mathfrak R_{\mathrm{res}}\) 负责离散 residue support；
  \item \(\mathfrak R_{\mathrm{amp}}\) 负责通过双切片恢复连续原子参数；
  \item \(\mathfrak R_{\mathrm{branch}}\) 负责 non-Abelian cubic singularity；
  \item \(\mathfrak R_{\mathrm{impl}}\) 负责真正的 multiplicative bookkeeping。
\end{enumerate}
等价地，一切最终只停留在“有限 cyclotomic 补丁 \(+\) prime-power Frobenius 塔 \(+\) 有限 Abelian residue ledger”上的 \(\pi\)-公式，都至多是值层算法，而不可能是 singularity-faithful 的 completion 公式。
\end{conjecture}

\endinput

\begin{theorem}[有限 cyclotomic--Hurwitz 阴影对 \texorpdfstring{$S_4$}{S4} 椭圆谱宿主的不完备性]\label{thm:conclusion-finite-cyclotomic-hurwitz-shadows-incomplete-for-s4-elliptic-host}
设 \(\mathfrak H_{\pi}\) 为项目的谱宿主，即由正规化
\[
y:E\to \PP^1
\]
与四次谱方程 \(\Pi(\lambda,y)\) 所承载的一参数椭圆覆盖。若一个不变量 \(\mathcal I\) 仅通过有限多个下列数据因子化：
\begin{enumerate}
  \item 根值单位平均
  \[
  \frac1m\sum_{j=0}^{m-1}a_n(\omega_m^j),
  \]
  \item 周期 Dirichlet 阴影
  \[
  Z_A(s)=d_m^{-s}\sum_{r=1}^{d_m}\kappa_A(r)\,\zeta\!\left(s,\frac{r}{d_m}\right),
  \]
\end{enumerate}
则 \(\mathcal I\) 不可能成为 \(\mathfrak H_{\pi}\) 的完备不变量。更精确地，它不能确定项目谱曲线的正规化 \(y:E\to \PP^1\) 及其四次谱方程 \(\Pi(\lambda,y)\) 的 \(\QQ(y)\)-同构类，尤其不能恢复
\[
\operatorname{Gal}_{\QQ(y)}(\Pi(\lambda,y))\cong S_4
\]
所携带的四层单值化几何。
\end{theorem}
\begin{proof}
由定理 \ref{thm:conclusion-realinput40-root-unity-ghost-complete-primitive-degenerate} 与定理 \ref{thm:conclusion-periodic-dirichlet-shadow-hurwitz-finite-closure}，这里允许的全部阴影数据只依赖于根值单位平均、周期计数与 Hurwitz 残数账本。它们本质上都是 sheet-symmetric 的标量压缩：前者只保留 ghost 层的整除梳齿，后者只保留有限周期向量及其平均值。

相反，文稿已证明项目谱宿主的正规化是一条固定椭圆曲线 \(E\)，而
\[
T_m(y)=\operatorname{Tr}_{\QQ(E)/\QQ(y)}(\lambda^m)
\]
形成真正的四层单值化迹族；相应四次谱方程在 \(\QQ(y)\) 上具有满 \(S_4\) 伽罗瓦群，其分歧与片交换结构正是宿主的核心几何数据。任何完备不变量若要恢复 \(\Pi(\lambda,y)\) 的 \(\QQ(y)\)-同构类，就必须保留这种非阿贝尔片排列信息。

然而，上述有限 cyclotomic--Hurwitz 阴影先天抹去片交换信息，只保留对所有 sheet 对称的标量影子。因此它们不可能完备恢复 \(\mathfrak H_{\pi}\) 的 \(S_4\) 单值化几何。证毕。
\end{proof}

\begin{theorem}[Puiseux--Hurwitz 二分律]\label{thm:conclusion-puiseux-hurwitz-dichotomy}
设 \(I(\alpha)\) 为项目 Lee--Yang 权重体系的速率函数。则当 \(\alpha\downarrow 0\) 与 \(\alpha\uparrow \frac12\) 时，分别有
\[
I(\alpha)=\log 2+2\alpha\log\alpha+(\log 8-2)\alpha+O\!\bigl(\alpha^2\log(1/\alpha)\bigr),
\]
\[
I\!\left(\frac12-\delta\right)
=
\log 2+4\delta\log\delta+(12\log 2-4)\delta
+O\!\bigl(\delta^2\log(1/\delta)\bigr),
\qquad \delta\downarrow 0.
\]
其中系数对
\[
(2,4)
\]
分别来自 Perron 分支在 \(y\to 0^+\) 与 \(y\to +\infty\) 处的 Puiseux 指数 \(1/2\) 与 \(1/4\)。相反，定理 \ref{thm:conclusion-periodic-dirichlet-shadow-hurwitz-finite-closure} 的任意有限 Hurwitz 阴影只携带 \(s=1\) 的简单极点及其残数账本。故 \((2,4)\) 不是 Hurwitz 残数账本内部的量，而是椭圆谱宿主的分支几何不变量。
\end{theorem}
\begin{proof}
端点展开及其 Puiseux 来源均已由文稿中的 Lee--Yang 端点分析严格给出。另一方面，定理 \ref{thm:conclusion-periodic-dirichlet-shadow-hurwitz-finite-closure} 表明：任何有限周期 Dirichlet 阴影在标量层面都分解为有限个 Hurwitz \(\zeta\) 函数之和，其奇性机制只表现为 \(s=1\) 处的简单极点及有限残数账本。

二者的解析型别因此严格不同：端点展开中的 \(\alpha\log\alpha\) 与 \(\delta\log\delta\) 来自代数分支的 Puiseux 指数，而 Hurwitz 阴影只记录 meromorphic 极点的残数。故端点系数对 \((2,4)\) 不可能在 Hurwitz 残数账本内部生成，只能来自谱宿主的 Puiseux 几何。证毕。
\end{proof}

\begin{conjecture}[三寄存器分裂原理]\label{conj:conclusion-ghost-hurwitz-sheet-three-register-splitting}
任一可精确嵌入项目的 \(\pi\)-通道，都应存在函子性分裂
\[
\mathfrak R_{\mathrm{ghost}}
\oplus
\mathfrak R_{\mathrm{Hurwitz}}
\oplus
\mathfrak R_{\mathrm{sheet}},
\]
其中
\[
\mathfrak R_{\mathrm{ghost}}
\quad\text{负责定理 \ref{thm:conclusion-realinput40-root-unity-ghost-complete-primitive-degenerate} 的 Dirichlet 梳齿层析，}
\]
\[
\mathfrak R_{\mathrm{Hurwitz}}
\quad\text{负责定理 \ref{thm:conclusion-periodic-dirichlet-shadow-hurwitz-finite-closure} 的有限周期残数账本，}
\]
\[
\mathfrak R_{\mathrm{sheet}}
\quad\text{负责定理 \ref{thm:conclusion-finite-cyclotomic-hurwitz-shadows-incomplete-for-s4-elliptic-host} 与定理 \ref{thm:conclusion-puiseux-hurwitz-dichotomy} 中的椭圆 \(S_4\) 单值化及端点 Puiseux 奇性。}
\]
并且，缺失第三寄存器的任何公式族都只能给出宿主的投影读出，而不可能与宿主本身同构；Machin-like 与 BBP-like 阴影应当都属于这一情形。
\end{conjecture}

\begin{theorem}[两条三重结构在覆盖范畴中绝不等价]\label{thm:conclusion-cf-two-trigonal-structures-never-equivalent}
设
\[
\phi_0,\phi_1:C_F\to \PP^1
\]
分别为文稿中由两个无穷远点诱导的第一、第二三重映射。则不存在
\[
(\sigma,\tau)\in \operatorname{Aut}(C_F)\times \operatorname{PGL}_2
\]
使
\[
\phi_1=\tau\circ \phi_0\circ \sigma .
\]
换言之，项目宿主不是“带一条三重解式”的对象，而是同时携带两条彼此不可相消的三重解式。
\end{theorem}
\begin{proof}
由定理 \ref{thm:fold-gauge-anomaly-first-trigonal-structure-golden-ratio} 与推论 \ref{cor:fold-gauge-anomaly-first-trigonal-monodromy-genus}，\(\phi_0\) 的分歧型为“两处全分歧 \(e=3\) 加六处简单分歧 \(e=2\)”，其 \(S_3\)-伽罗瓦闭包亏格为
\[
g=8.
\]
而由定理 \ref{thm:fold-gauge-anomaly-second-trigonal-structure-discriminant} 与推论 \ref{cor:fold-gauge-anomaly-second-trigonal-monodromy-genus}，\(\phi_1\) 的全部十个分歧参数皆为简单分歧，且其 \(S_3\)-伽罗瓦闭包亏格为
\[
g=10.
\]
覆盖在 \(\operatorname{Aut}(C_F)\times \operatorname{PGL}_2\) 作用下，分歧指数多重集与伽罗瓦闭包亏格皆保持不变。由于这两组不变量不同，\(\phi_0\) 与 \(\phi_1\) 不可能等价。证毕。
\end{proof}

\begin{theorem}[有限个三次纤维绝不生成第二分歧九次域]\label{thm:conclusion-second-branch-ninefold-field-not-generated-by-finitely-many-cubic-fibers}
记 \(L_9/\QQ\) 为第二三重结构九次判别多项式 \(P_9(t)\) 的分裂域。对每个满足
\[
t\ge 5,\qquad t\not\equiv 1\pmod 3
\]
的整数 \(t\)，令 \(M_t/\QQ\) 为三次专化多项式 \(h_t(u)\) 的分裂域。则对任意有限集合
\[
t_1,\dots,t_r,
\]
都有
\[
L_9\not\subset M_{t_1}\cdots M_{t_r}.
\]
\end{theorem}
\leanverified{paper\_conclusion\_second\_branch\_ninefold\_field\_not\_generated\_by\_finitely\_many\_cubic\_fibers}
\begin{proof}
由定理 \ref{thm:fold-gauge-anomaly-P9-galois-discriminant}，
\[
\Gal(L_9/\QQ)\cong S_9,
\]
故 \(L_9/\QQ\) 非可解。另一方面，由定理 \ref{thm:fold-gauge-anomaly-second-trigonal-ht-specialization-s3-disjoint}，对每个上述 \(t\) 都有
\[
\Gal(M_t/\QQ)\cong S_3.
\]
于是有限复合域
\[
M_{t_1}\cdots M_{t_r}
\]
的伽罗瓦闭包群嵌入到有限直积 \(S_3^r\) 中，从而为可解群。非可解扩张 \(L_9/\QQ\) 不可能嵌入一个可解扩张之内，故
\[
L_9\not\subset M_{t_1}\cdots M_{t_r}.
\]
证毕。
\end{proof}

\begin{corollary}[判别式平方专化仅为有限异常层]\label{cor:conclusion-second-trigonal-square-discriminant-only-finite-exceptional-layer}
使得
\[
\Disc_u(h_t)\in \QQ^{\times 2}
\]
成立的有理参数 \(t\in \QQ\) 仅有有限多个。故循环三次专化不构成 generic 现象，而只形成一个有限异常层。
\end{corollary}
\begin{proof}
这正是定理 \ref{thm:fold-gauge-anomaly-P9-discriminant-square-finiteness} 的直接重述。该定理把平方判别式条件等价转写为属 \(4\) 的超椭圆曲线
\[
\mathcal H:\quad y^2=-(t-1)P_9(t)
\]
上的有理点条件，并由 Faltings 定理推出 \(\mathcal H(\QQ)\) 有限，故满足判别式平方条件的 \(t\) 也只有有限多个。证毕。
\end{proof}

\begin{theorem}[单寄存器 Chebotarev 证书不可能完备]\label{thm:conclusion-single-register-chebotarev-cannot-be-complete-for-double-branch-arithmetic}
记 \(L_{10}\) 为第一分歧十次域，\(L_9\) 为第二分歧九次域。则不存在任何仅因子化于
\[
\Gal(L_{10}/\QQ)\cong S_{10}
\quad\text{或仅因子化于}\quad
\Gal(L_9/\QQ)\cong S_9
\]
之一的 Frobenius 证书，能够完备决定 projective host 的双分歧算术。
\end{theorem}
\begin{proof}
由定理 \ref{thm:fold-gauge-anomaly-P10-P9-linear-disjointness}，
\[
L_{10}\cap L_9=\QQ,
\qquad
\Gal(L_{10}L_9/\QQ)\cong S_{10}\times S_9.
\]
Chebotarev 密度定理在直积群上给出乘积分布。于是对任意真共轭类子集
\[
C_\lambda\subset S_{10},
\qquad
C_\mu\subset S_9,
\]
两类素数集合
\[
\{p:\ \Frob_p\in C_\lambda\times C_\mu\},
\qquad
\{p:\ \Frob_p\in C_\lambda\times (S_9\setminus C_\mu)\}
\]
只要边际类非空，就都具有正自然密度。故在固定第一寄存器的条件下，第二寄存器仍存在正密度摆动；反之亦然。于是任何只读取一个投影坐标的证书，必在正密度集合上混同互异的双分歧算术。

特别地，由推论 \ref{cor:fold-gauge-anomaly-P10-P9-chebotarev-independence}，
\[
\delta\bigl(\text{\(P_{10}\bmod p\) 不可约且 \(P_9\bmod p\) 不可约}\bigr)=\frac1{90},
\]
这已是双寄存器独立性的显式见证。证毕。
\end{proof}

\begin{theorem}[Lee--Yang 局部分歧零包的二值分类律]\label{thm:conclusion-leyang-local-doublezero-binary-classification}
\leanverified{paper\_conclusion\_leyang\_local\_doublezero\_binary\_classification}
设 \(u_\ast\) 为主稿第 \(6.219\)–\(6.222\) 节中的一个非平凡分歧点，并考虑靠近 \(u_\ast\) 的局部零包
\[
u_{m,1},u_{m,2},\dots
\]
按与 \(u_\ast\) 的距离递增排序。则存在两且仅两种普适类型：
\begin{enumerate}
  \item \(\sinh\)-型：
  \[
  \frac{|u_{m,2}-u_\ast|}{|u_{m,1}-u_\ast|}\longrightarrow 4;
  \]
  \item \(\cosh\)-型：
  \[
  \frac{|u_{m,2}-u_\ast|}{|u_{m,1}-u_\ast|}\longrightarrow 9.
  \]
\end{enumerate}
因此，仅凭最邻近两枚局部零点的尺度比，便可判别分歧点所属的普适核型。
\end{theorem}
\begin{proof}
主稿已经证明：局部归一化极限只能落入
\[
2C_1\frac{\sinh w}{w}
\qquad\text{或}\qquad
2C_0\cosh w
\]
两类，其中前者的极限零点为 \(w=i\pi k\)（\(k\in\ZZ\setminus\{0\}\)），后者的极限零点为 \(w=i\pi(k+\tfrac12)\)。再由局部反变换
\[
u=u_\ast+\frac{w^2}{b_\ast m^2}
\]
可知
\[
u_{m,j}-u_\ast
=
-\frac{\pi^2\eta_j}{b_\ast m^2}+o(m^{-2}),
\]
其中
\[
\eta_j=
\begin{cases}
j^2,& \sinh\text{-型},\\[2mm]
\left(j-\tfrac12\right)^2,& \cosh\text{-型}.
\end{cases}
\]
于是
\[
\frac{|u_{m,2}-u_\ast|}{|u_{m,1}-u_\ast|}
\longrightarrow
\frac{\eta_2}{\eta_1}
=
\begin{cases}
4,& \sinh\text{-型},\\
9,& \cosh\text{-型}.
\end{cases}
\]
证毕。
\end{proof}

\begin{theorem}[两零点 microlocal tomography 对 \((u_\ast,b_\ast)\) 的重构公式]\label{thm:conclusion-leyang-twozero-microlocal-tomography}
沿用定理 \ref{thm:conclusion-leyang-local-doublezero-binary-classification} 的记号。若已由二值分类律判定普适类型，并令
\[
(\eta_1,\eta_2)=
\begin{cases}
(1,4),& \sinh\text{-型},\\[2mm]
(\tfrac14,\tfrac94),& \cosh\text{-型},
\end{cases}
\]
则由同一尺度 \(m\) 上最邻近的两枚局部零点，可重构
\[
b_\ast
=
-\,\frac{\pi^2(\eta_2-\eta_1)}{m^2\,(u_{m,2}-u_{m,1})}+o(1),
\]
以及
\[
u_\ast=
\frac{\eta_2u_{m,1}-\eta_1u_{m,2}}{\eta_2-\eta_1}+o(m^{-2}).
\]
因此，分歧点位置与 Puiseux 斜率平方不变量并不需要全局谱曲线；它们已经被双零点局部包完全编码。
\end{theorem}
\begin{proof}
由定理 \ref{thm:conclusion-leyang-local-doublezero-binary-classification} 的渐近式，
\[
u_{m,j}
=
u_\ast-\frac{\pi^2\eta_j}{b_\ast m^2}+o(m^{-2})
\qquad (j=1,2).
\]
两式相减得到
\[
u_{m,2}-u_{m,1}
=
-\frac{\pi^2(\eta_2-\eta_1)}{b_\ast m^2}+o(m^{-2}),
\]
故
\[
b_\ast
=
-\,\frac{\pi^2(\eta_2-\eta_1)}{m^2\,(u_{m,2}-u_{m,1})}+o(1).
\]
再以 \(\eta_2\) 倍第一式减去 \(\eta_1\) 倍第二式，主项中的 \(b_\ast\)-项相消，得到
\[
(\eta_2-\eta_1)u_\ast
=
\eta_2u_{m,1}-\eta_1u_{m,2}+o(m^{-2}),
\]
从而
\[
u_\ast=
\frac{\eta_2u_{m,1}-\eta_1u_{m,2}}{\eta_2-\eta_1}+o(m^{-2}).
\]
证毕。
\end{proof}

\begin{corollary}[局部零包决定整个十次数域]\label{cor:conclusion-leyang-local-doublezero-recovers-decic-field}
沿用定理 \ref{thm:conclusion-leyang-twozero-microlocal-tomography} 的设定。由主稿第 \(6.222\) 节，
\[
\QQ(u_\ast)=\QQ(\mu_\ast)=\QQ(b_\ast),
\]
且该公共数域为 \(\QQ\) 的十次扩张。再由方向公式
\[
\operatorname{dir}(u_\ast)=-\frac{1}{b_\ast},
\]
进一步有
\[
\QQ(\operatorname{dir}(u_\ast))
=
\QQ(u_\ast)
=
\QQ(\mu_\ast)
=
\QQ(b_\ast).
\]
因此，一个分歧点附近的双零点局部包，不仅决定局部几何常数，而且决定整个十次分歧数域的同构类。
\end{corollary}
\begin{proof}
由定理 \ref{thm:conclusion-leyang-twozero-microlocal-tomography}，双零点局部包已恢复 \(u_\ast\) 与 \(b_\ast\)。主稿又已证明
\[
\QQ(u_\ast)=\QQ(\mu_\ast)=\QQ(b_\ast),
\]
且这一公共数域的次数为 \(10\)。由于
\[
\operatorname{dir}(u_\ast)=-1/b_\ast
\]
与 \(b_\ast\) 互为有理函数，故
\[
\QQ(\operatorname{dir}(u_\ast))=\QQ(b_\ast).
\]
综上，双零点局部包已决定整个十次数域。证毕。
\end{proof}

\begin{theorem}[Lee--Yang 边界并合核的二核普适性与十次数域刚性]\label{thm:conclusion-leyang-boundary-coalescence-binary-kernel-decic-rigidity}
设 \(u_\ast\) 为主稿第 \(6.222\)–\(6.225\) 节中的一个非平凡主导谱支并合点，并沿用其中的局部归一化记号 \(\widetilde M_m\)、分支值 \(\mu_\ast\) 与方向不变量 \(b_\ast\)。则恰有两种且仅有两种局部普适核：
\[
\widetilde M_m\!\left(u_\ast+\frac{w^2}{b_\ast m^2}\right)
\longrightarrow
2C_1\frac{\sinh w}{w},
\]
或
\[
\widetilde M_m\!\left(u_\ast+\frac{w^2}{b_\ast m^2}\right)
\longrightarrow
2C_0\cosh w.
\]
相应地，局部零点满足统一二次晶格律
\[
u_{m,k}
=
u_\ast-\frac{\pi^2(k+\varepsilon)^2}{b_\ast m^2}
+O\!\left(\frac{k}{m^3}\right),
\qquad
\varepsilon\in\left\{0,\frac12\right\}.
\]
此外，
\[
\QQ(u_\ast)=\QQ(\mu_\ast)=\QQ(b_\ast),
\qquad
[\QQ(b_\ast):\QQ]=10.
\]
因此，非平凡 Lee--Yang 边界并合不存在第三种局部普适核，而其方向不变量也不是连续模参数，而是严格锁定在一个十次数域中的代数点。
\end{theorem}
\begin{proof}
主稿第 \(6.222\) 节已经证明：在给定归一化下，局部极限只能落入
\[
2C_1\frac{\sinh w}{w}
\qquad\text{或}\qquad
2C_0\cosh w
\]
两类，且不存在第三种情形。定理 \ref{thm:conclusion-leyang-local-doublezero-binary-classification} 只是把这两类核进一步压缩为最近两枚局部零点尺度比的二值判据。

随后，主稿第 \(6.223\) 节给出零点的二次晶格展开；其中 \(\varepsilon=0\) 与 \(\varepsilon=\tfrac12\) 分别对应 \(\sinh\)-型与 \(\cosh\)-型。最后，推论 \ref{cor:conclusion-leyang-local-doublezero-recovers-decic-field} 已经说明
\[
\QQ(u_\ast)=\QQ(\mu_\ast)=\QQ(b_\ast)=\QQ(\operatorname{dir}(u_\ast)),
\]
且该公共数域的次数为 \(10\)。三者合并即得结论。证毕。
\end{proof}

\begin{theorem}[Lee--Yang 边缘对数导数的 Dirichlet--resolvent 极限]\label{thm:conclusion-leyang-edge-log-derivative-dirichlet-resolvent-limit}
\leanverified{paper\_conclusion\_leyang\_edge\_log\_derivative\_dirichlet\_resolvent\_limit}
沿用结论 \ref{con:xi-terminal-zm-double-root-diffusive-sinc-limit} 的多项式族 \(Z_m(y)\)，并定义边缘缩放对数导数
\[
R_m(u):=
-\frac{1}{m^2}\,\partial_y\log Z_m(y)\Big|_{y=-u/m^2}.
\]
则对任意紧集
\[
K\Subset \CC\setminus \{4\pi^2j^2:\ j\ge 1\}
\]
成立局部一致极限
\[
R_m(u)\longrightarrow \frac{S'(u)}{S(u)}
=-\sum_{j\ge 1}\frac{1}{4\pi^2j^2-u},
\qquad u\in K,
\]
其中
\[
S(u)=\operatorname{sinc}\!\left(\frac{\sqrt u}{2}\right)
=\prod_{j\ge 1}\left(1-\frac{u}{4\pi^2j^2}\right).
\]
\end{theorem}
\begin{proof}
由结论 \ref{con:xi-terminal-zm-double-root-diffusive-sinc-limit}，
\[
\frac{Z_m(-u/m^2)}{N_m}\longrightarrow S(u)
\]
在任意紧集上一致成立，而 \(K\) 避开 \(S\) 的零点集合 \(\{4\pi^2j^2\}_{j\ge1}\)。
因此对充分大的 \(m\)，\(Z_m(-u/m^2)\) 在 \(K\) 的某个邻域内无零点，故可取解析对数并应用局部一致收敛的解析函数对数导数稳定性，得到
\[
\partial_u\log Z_m(-u/m^2)\longrightarrow \partial_u\log S(u)
\]
于 \(K\) 上局部一致。
又
\[
\partial_u\log Z_m\!\left(-\frac{u}{m^2}\right)
=
-\frac{1}{m^2}\,\partial_y\log Z_m(y)\Big|_{y=-u/m^2}
=R_m(u),
\]
故 \(R_m(u)\to S'(u)/S(u)\)。
最后，结论 \ref{con:xi-terminal-zm-diffusive-edge-weyl-spectral-determinant} 已给出 Weierstrass 乘积
\[
S(u)=\prod_{j\ge 1}\left(1-\frac{u}{4\pi^2j^2}\right),
\]
逐项取对数导数即得
\[
\frac{S'(u)}{S(u)}=-\sum_{j\ge 1}\frac{1}{4\pi^2j^2-u}.
\]
证毕。
\end{proof}

\begin{theorem}[Lee--Yang 边缘零点测度的 Dirichlet 时钟收敛]\label{thm:conclusion-leyang-edge-zero-measure-dirichlet-clock-convergence}
沿用结论 \ref{con:xi-terminal-zm-double-root-diffusive-sinc-limit} 的记号，记
\[
\rho_{m,1}>\rho_{m,2}>\cdots
\]
为 \(Z_m\) 在 \((-\infty,0)\) 上由右向左排列的零点，并定义缩放原子测度
\[
\mu_m:=\sum_{j\ge 1}\delta_{m^2(-\rho_{m,j})}.
\]
则在 \((0,\infty)\) 上有 vague 收敛
\[
\mu_m \xrightarrow{\mathrm{vague}} \mu_D:=\sum_{j\ge 1}\delta_{4\pi^2j^2}.
\]
等价地，对任意 \(\varphi\in C_c((0,\infty))\)，有
\[
\lim_{m\to\infty}\sum_{j\ge1}\varphi\!\bigl(m^2(-\rho_{m,j})\bigr)
=
\sum_{j\ge1}\varphi(4\pi^2j^2).
\]
\end{theorem}
\begin{proof}
由结论 \ref{con:xi-terminal-zm-double-root-diffusive-sinc-limit}，对每个固定 \(j\ge1\) 都有
\[
m^2(-\rho_{m,j})\longrightarrow 4\pi^2j^2.
\]
另一方面，结论 \ref{con:xi-terminal-zm-diffusive-edge-weyl-spectral-determinant} 已给出局部计数律：对任意固定 \(U>0\)，
\[
\#\{j:\ m^2(-\rho_{m,j})\le U\}
=
\left\lfloor \sqrt{\frac{U}{4\pi^2}}\right\rfloor+o(1).
\]
因此任意紧支撑测试函数 \(\varphi\) 只会看到有限多个原子，且不会发生质量逃逸。
逐个原子的收敛与局部总质量控制合并，即得所述 vague 收敛。证毕。
\end{proof}

\begin{theorem}[Lee--Yang 四次族的解析边缘与算术纤维双普适分裂]\label{thm:conclusion-leyang-analytic-edge-arithmetic-fiber-dual-universality-split}
同一四次族 \(\Pi(\lambda,y)=0\) 同时呈现两种彼此不可约化的普适极限对象：
\begin{enumerate}
\item 在边缘缩放 \(y=-u/m^2\) 下，出现由
\[
\mu_D=\sum_{j\ge1}\delta_{4\pi^2j^2}
\]
支撑的确定性 Dirichlet 时钟，其一点评价由 meromorphic resolvent
\[
u\longmapsto -\sum_{j\ge1}\frac{1}{4\pi^2j^2-u}
\]
完全编码；
\item 在良素未分歧参数上，纤维根数 \(N_p(y_0)\) 的低阶统计呈现 quartic Poisson 伪装：其概率生成函数满足
\[
\frac{1}{|U_p^{\mathrm{aff}}|}\sum_{y_0\in U_p^{\mathrm{aff}}} z^{N_p(y_0)}
=
\frac38+\frac13 z+\frac14 z^2+\frac1{24}z^4+O_z(p^{-1/2}),
\]
从而一切四次以下多项式观测都与 \(\mathrm{Pois}(1)\) 一致到 \(O(p^{-1/2})\)，而五阶下降阶乘是首个精确可见破缺阶。
\end{enumerate}
因此，该 Lee--Yang 四次族在解析边缘侧给出“无限极点、即时显影”的确定性时钟，而在算术纤维侧给出“前四阶全隐、第五阶才破”的随机 Poisson 伪装；二者共享同一代数方程，却属于两种不同的可识别机制。
\end{theorem}
\begin{proof}
第一部分正是定理 \ref{thm:conclusion-leyang-edge-log-derivative-dirichlet-resolvent-limit} 与定理 \ref{thm:conclusion-leyang-edge-zero-measure-dirichlet-clock-convergence} 的内容。
第二部分由定理 \ref{thm:conclusion-leyang-unramified-fiber-count-quartic-poisson-cloak} 与推论 \ref{cor:conclusion-leyang-fifth-falling-factorial-first-visible-break} 给出。
两侧的极限对象分别由 meromorphic resolvent 与有限阶多项式统计承载，且第二侧直到五阶才出现首个不可消去破缺，故它们不能压缩为同一种有限阶 one-point 观测机制。证毕。
\end{proof}

\begin{conjecture}[tame \texorpdfstring{$S_d$}{Sd} 覆盖的 degree-\texorpdfstring{$d$}{d} Poisson 伪装原理]\label{conj:conclusion-leyang-degree-d-poisson-cloak-principle}
设
\[
\pi:C\to \PP^1
\]
是一个 tame good reduction 的 \(S_d\)-generic 有理覆盖。
对良素 \(p\) 与未分歧参数 \(y_0\) 定义纤维根数
\[
N_p(y_0):=\#\pi^{-1}(y_0)(\FF_p).
\]
则其 good-prime 极限分布应满足：
\[
\EE\bigl[(N_p)_k\bigr]=1+O(p^{-1/2})
\qquad (1\le k\le d),
\]
而
\[
(N_p)_{d+1}\equiv 0
\]
逐点成立。
等价地，其概率生成函数在 \(z=1+t\) 处应满足
\[
\EE\bigl[(1+t)^{N_p}\bigr]
=
\sum_{k=0}^{d}\frac{t^k}{k!}+O(p^{-1/2}).
\]
换言之，good-prime 纤维根数在低阶统计上普遍伪装成参数 \(1\) 的 Poisson 律，而这种伪装的长度恰由覆盖次数 \(d\) 控制；在本节所处理的 Lee--Yang 四次族中，这一原理的 \(d=4\) 情形已被完全验证。
\end{conjecture}

\endinput


综上，本节相关结构并不集中于单一数域或单一 Fourier 机制，而是同时由 \(K_{\mathrm{LY}}=\QQ(\sqrt{-111})\) 所控制的二次分辨层、\(K_{\mathrm{br}}=\QQ(u)\) 所承载的非正规三次层、\(\QQ(\sqrt5)\) 所控制的局部正则二次层、\(\QQ(B)\) 所引入的六次振幅层、\(\TT_2\) 上不可 Hilbert--Schmidt 的奇异卷积层，以及 \(B_\infty\) 所定义的零温四态地基动力学共同承载。尤其在 Lee--Yang 奇性坐标上，全局正规化、阿贝尔化与 \(5\)-进完成三者彼此不交换：全局最小正规闭包被迫长成 \(C_2\times S_3\) 型对象，其最大阿贝尔阴影却只剩双二次层，而所有 \(5\)-进局部完成又同步退化为循环局部纤维。与此同时，single-atom residue 的连续参数恢复并不发生在同一算术层：单切片始终保留一维乘法轨道，有限 cyclotomic 补丁对正 twisted 指数完全透明，因而 \((m,E)\) 的 faithful 读出至少需要双切片；而 prime-power 完成化塔在 \(\bmod p\) 上又整体塌缩为 \(\FF_p[s^p]\) 的单生成 Frobenius 轨道。进一步说，这种不交换性与塌缩性合并后，会迫使 residue support、双切片幅度、non-Abelian singular branch 与真正的 multiplicative bookkeeping 分裂到彼此不可消去的层上：根值单位平均只负责 ghost 梳齿层析，有限周期 Dirichlet 阴影只负责 Hurwitz 残数账本，而真正的椭圆 \(S_4\) 单值化与端点 Puiseux 指数则被迫落在独立的 sheet 寄存器中；若再要求 faithful 的乘法实现，则还必须额外外置实现层。在分歧点的 microlocal 层面，这一 sheet 结构又已被双零点局部包刚性编码，并直接恢复 \((u_\ast,b_\ast)\) 以及它们所张成的十次数域。于是相关非阿贝尔性既不诞生于单一局部坏素，也不诞生于单一正则 germ，而是诞生于正则黄金二次层与奇性三次层在全局拼接时所形成的最小不可交换闭包。

\endinput
